% Exact current Paper I cited input definitions and statements for II GI.
% No II GI/GK/GN target proof or old computation is included.

% standard affine and cusp data, full statement; exact current I lines 8916--8940
\begin{proposition}[Classification with arbitrary affine and cusp data]\hypertarget{tag-01QB}{}\label{fam:src:B.2.1}\label{fam:st:B.2.1}
\StatementRecord{fam:src:B.2.1}{proved}{01QB}{\SeriesRef{fam:src:A.1.1}, \SeriesRef{fam:src:A.3.1}; \cite[Section~6.4, pp.~39--40]{AC2026}}

Let
\[
 w_1=r_1\varepsilon+k_1d,\qquad w_2=r_2\varepsilon'+k_2d,
 \qquad 3\nmid r_1,\quad 2\nmid r_2,
\]
and use an arbitrary holomorphic cusp collar section $\sigma$ in the construction of \cite[Section~6.4, pp.~39--40]{AC2026}. Write $l_0=\gamma(c(\sigma))$ in its positive cusp-winding convention, and set $p=12l_0-4r_1-3r_2$. Let primes denote a second such datum with the same period solution and cusp fan. Then
\[
 X(c;w_1,w_2,\sigma)\simeq X(c';w'_1,w'_2,\sigma')
\]
holds exactly for the following arithmetic conditions:
\[
 n=c'-c\in\mathbb Z,\qquad
 p'=a p\text{ for some }a\in\{1,-1\},\qquad
 n\equiv k_2-k'_2\pmod2.
\]
All maps are over $B$. A map with the indicated $a$ has fibre-linear part $aI_2$. In particular the complete invariants are
\[
 \bigl(|p|,\ e^{-\pi i(c+k_2)}\bigr),\qquad
 |p|\ge1,\quad\gcd(|p|,6)=1.
\]
Every value of this positive integer occurs. The full collar section enters the construction, but its additional analytic part and its second integral winding contribute no further invariant.
\end{proposition}

% higher multiplicities, full statement; exact current I lines 9047--9057
\begin{proposition}[Classification for higher multiplicities]\hypertarget{tag-01QE}{}\label{fam:src:B.3.1}\label{fam:st:B.3.1}
\StatementRecord{fam:src:B.3.1}{proved}{01QE}{\SeriesRef{fam:src:A.1.2}, \SeriesRef{fam:src:A.2.1}, \SeriesRef{fam:src:B.2.1}}

Let $k,l$ be positive integers with $\gcd(3k,4l)=1$, and put $m_1=3k,m_2=4l$. Consider Viaclovsky's sphere members expressed in the same clockwise marking, with $w_j=r_j\varepsilon_j+s_jd$, $\gcd(r_j,m_j)=1$, cusp winding $l_0$, and $p=m_1m_2l_0-m_2r_1-m_1r_2=\pm1$, where $\varepsilon_1=\varepsilon,\varepsilon_2=\varepsilon'$. Two such members are biholomorphic precisely when their multiplicity pair agrees and, for some $a\in\{1,-1\}$ and integer $n=c'-c$,
\[
 p'=ap,\quad r'_j\equiv ar_j\pmod{m_j},\quad
 nr'_1\equiv as_1-s'_1\pmod{k},\quad
 nr'_2\equiv as_2-s'_2\pmod{2l}.
\]
The congruences on $r'_j$ already follow from $p'=ap$. For fixed discrete vectors and cusp collar the identification subgroup is $2kl\mathbb Z$. For fixed admissible $u$ and this multiplicity pair, all the stated sphere translation data give exactly $2kl$ isomorphism classes, permuted in one cycle by a circuit of $u$ about zero.
\end{proposition}

% source-coordinate definition and integral dictionary statement; exact current I lines 9188--9234
\subsection{Higher logarithmic transformations}\hypertarget{tag-02ZE}{}\label{push:section:E.1}
\begin{definition}[Higher multiplicities and source coordinates]
\hypertarget{tag-IP10}{}\label{I:def:higher-multiplicity-coordinates}
\StatementRecord{I:def:higher-multiplicity-coordinates}{proved}{IP10}{}
Put $m=3k$, $n=4l$ with $\gcd(m,n)=1$.
Viaclovsky's indices $(0,1)$ correspond to the finite indices $(1,2)$
of \cref{def:common-conventions}. Lattice coordinates are integral
homology coordinates, and matrices act on columns.
\end{definition}

\begin{proposition}[The integral dictionary]\hypertarget{tag-032R}{}\label{push:src:E.1.1}
\StatementRecord{push:src:E.1.1}{proved}{032R}{\cite[Section~2, equations~(2.1)--(2.2), p.~3, and (2.7), (2.11)--(2.12), p.~4]{ViaclovskyLog2026}}

Let $(\delta,a,b,d)$ be the marking in
\cite[Section~2, equations~(2.1)--(2.2), p.~3, and (2.7), (2.11)--(2.12), p.~4]{ViaclovskyLog2026},
and put $u_0=d+2a+2b$, $u_1=d+3a+3b$.
The homomorphism from the marking of \cref{def:common-conventions} to this marking is
\[
 P_V=\begin{pmatrix}
 -2&1&0&1\\-6&0&-1&0\\0&-1&0&0\\-1&0&0&0
 \end{pmatrix},\qquad \det P_V=-1.
\]
It intertwines the two finite monodromies and satisfies
\[
 P_V\widehat\delta=\delta,\qquad
 P_V\varepsilon=-u_0,\qquad
 P_V\varepsilon'=\delta-u_1.
\]
With the counterclockwise deck generators of that paper, its data
$m_0=3k$, $m_1=4l$ and
$a_i=[v_i/m_i]$ with $v_i=\rho_i\delta+N_i u_i$ correspond to
\[
 w_1=N_0\varepsilon-\rho_0\widehat\delta,
 \qquad
 w_2=N_1\varepsilon'-(\rho_1+N_1)\widehat\delta.
\]
For the section-preserving cusp comparison used there, $l_0=0$ and
$p=-(nN_0+mN_1)$.  Thus its sphere condition is
$b:=nN_0+mN_1=\pm1$.
The brackets in $a_i$ mean the torsion point in the smooth
root torus, using its section-preserving origin and the
transported period vector.  Replacing $v_i$ by $v_i+m_i\lambda$, $\lambda\in\Lambda$,
leaves the class $[v_i/m_i]$ unchanged, while its chosen lift fixes the integral collar defect.
\end{proposition}


\begin{proof}

% residual cusp-defect definition; exact current I lines 9256--9270
\begin{definition}[The residual cusp defect]
\hypertarget{tag-IP11}{}\label{I:def:residual-cusp-defect}
\StatementRecord{I:def:residual-cusp-defect}{proved}{IP11}{}
For a general cusp collar $\sigma$, write $c(\sigma)$ for
its residual translation defect in
$\overline\Lambda=\Lambda/(\mathbb Z\widehat w+
\mathbb Z\widehat\delta)$, using the convention
$\sigma(g_0z)-R_0\sigma(z)=\Pi_c(g_0z)\ell(\sigma)$
and $c(\sigma)=[\ell(\sigma)]$.
Set $l_0=\gamma(c(\sigma))$.  In the marked coordinates
$w_j=r_j\varepsilon_j+s_j\widehat\delta$, the integer is
$p=mn l_0-nr_1-mr_2$.
The section-preserving source data of \ref{push:src:E.1.1} have $c(\sigma)=0$.
\end{definition}


% current cited input tangent theorem statement; exact current I lines 31457--31494
\begin{theorem}[Tangent cohomology, direct images and miniversality]\hypertarget{tag-032U}{}\label{push:src:E.2.1}
\hypertarget{tag-01RH}{}\label{fam:src:A.5.2}\label{fam:st:A.5.2}\label{inv:D1.3}\label{inv:D2.1}\label{inv:D3.1}\label{cdp:st:E9}
\hypertarget{tag-02YO}{}\label{night:P5.3}\label{stmt:P5.3}
\hypertarget{tag-02YP}{}\label{night:P5.4}\label{stmt:P5.4}
\StatementRecord{push:src:E.2.1}{proved}{032U}{\cref{push:section:E.1,push:src:E.3.1,inv:D1.1,inv:D1.2}}
Let $k,l$ be positive integers with $\gcd(3k,4l)=1$, and let $X$
be a sphere member with multiplicities $(3k,4l)$ and the admissible
logarithmic data of \cref{push:section:E.1}. Put $z=kl$ and
$d=3k+4l$. On $\mathfrak B_{k,l}=\mathbb P(3k,4l)$, the vertical
root bundle $E$, the root lift $\bar f$ and its direct images satisfy
\[
 \begin{split}
 E&=\mathcal O\oplus\mathcal O(-z),\\
 R^q\bar f_*\mathcal O_X&=
 \mathcal O,\quad\mathcal O\oplus\mathcal O(-z),\quad\mathcal O(-z),\\
 K_X&=\bar f^*\mathcal O(z-d).
 \end{split}
\]
The tangent cohomology is
$\bigl(h^0,h^1,h^2,h^3\bigr)(X,T_X)=(1,1,1,0)$.
For the transition-arrow convention
$\kappa(\Phi)=(\partial_c\Phi)\Phi^{-1}$, the Kodaira--Spencer
class of $\partial_c$ is $e\otimes[\gamma]$, and the class of
$\partial_u$ is $-(2\pi i u)^{-1}e\otimes[\gamma]$.
Both classes are nonzero. The primary Kuranishi bracket vanishes,
the Kuranishi germ is a smooth reduced disc, and the fixed-data
period family is locally miniversal.

For the standard $(3,4)$ data, $e=\xi$ and the coarse direct
images are
\begin{equation}
 f_*T_X=\mathcal O\oplus\mathcal O(-1),\qquad
 R^1f_*T_X=\mathcal O\oplus\mathcal O(-1)^2,\qquad
 R^2f_*T_X=\mathcal O.
\label{pub:I:eq:D-4}
\end{equation}
\end{theorem}

